\section{High ranks at a fixed all-real time}\label{sec:comparison}

Fix \(T>0\) at which the full divisor is known to be real. Set
\[
 F(z)=H(T,z)/Z(T),\qquad e_m=\frac{\E_T U^{2m}}{(2m)!},\qquad
 G(v)=\sum_{m\ge0}(-1)^m e_m v^m .
\]
Thus \(F(z)=G(z^2)\). We compare two uses of these same coefficients. The
first estimates an individual factor from symmetric functions. The second
counts polynomial zeros on circles where the full field has a margin
\cite{Founding07,EZHighRank}.

\subsection{The complete product and fixed-time spacing}

For \(u\ge0\), the source in \eqref{eq:field} is positive and is bounded
above by \(C e^{9u-\pi e^{4u}}\). Indeed, after factoring out the first
exponential, the sum of the upper bounds \(2\pi^2 n^4 e^{-\pi(n^2-1)}\)
is finite. For \(|z|\le r\), use \(|\cos(zu)|\le e^{ru}\).
Maximizing \(ru-\pi e^{4u}/2\) on the half line bounds it above by
\(O(r\log(2+r))\). The integral with the remaining factor
\(e^{Tu^2+9u-\pi e^{4u}/2}\) is finite. Hence the logarithm of the
maximum modulus is \(O_T(1+r\log(2+r))\), and Jensen's formula gives
\(O_T(r\log(2+r))\) zeros in the disk of radius \(r\) for \(r\ge2\).
In particular \(\sum_i x_i^{-2}<\infty\), counting positive zeros with
multiplicity. The entire function has order at most one, so Hadamard
factorization pairs the genus-one factors at \(\pm x_i\). Evenness removes
the linear exponential, and \(F(0)=1\) fixes the constant. Therefore
\begin{equation}\label{eq:all-real-product}
 F(z)=\prod_{i\ge1}(1-b_i z^2),\qquad
 G(v)=\prod_{i\ge1}(1-b_i v),\qquad b_i=x_i^{-2}.
\end{equation}
The products converge locally uniformly. In particular \(e_m\) is the
\(m\)-th elementary symmetric sum of the complete sequence \((b_i)\).

Ki, Kim and Lee's Theorems 1.3--1.4 give the simple real tail and its
asymptotic spacing at each fixed positive time \cite{KiKimLee2009}.
Their normalization is \(\Xi_\mu(w)=8H(4\mu,2w)\), so in physical
coordinates \(\Delta_i\log x_i\to4\pi\), where
\(\Delta_i=x_{i+1}-x_i\). The all-real assumption identifies this tail
with the positive product ranks up to a finite shift. Applying the
mean-value theorem to \(x\log x-x\) across each such gap gives increments
converging to \(4\pi\), since \(\Delta_i/x_i\to0\).
Telescoping and inverting gives
\begin{equation}\label{eq:high-rank-geometry}
 x_i\sim\frac{4\pi i}{\log i},\qquad
 \Delta_i\sim\frac{4\pi}{\log i},\qquad
 \sum_{i>j}b_i\ll_T\frac{\log^2(2j)}{j}.
\end{equation}
The tail-sum bound follows by comparison with
\(\int_j^\infty \log^2(2v)\,\dd v/v^2\). The onset of this spacing
regime at the chosen time is unevaluated.

\subsection{Recovering individual factors}

Let \(B_j=\prod_{i\le j} b_i\), \(B_0=1\), and let \(E_{j,h}\) be
the \(j\)-th elementary symmetric sum of the sequence \((b_i^h)_i\)
for integer \(h\ge1\). In particular \(E_{0,h}=1\). The logarithm
of \eqref{eq:all-real-product} near zero and Newton's identities give
\[
 p_n=-n[v^n]\log G(v)=\sum_i b_i^n,qquad
 E_{j,h}=\frac1j\sum_{n=1}^j(-1)^{n-1}p_{nh} E_{j-n,h} .
\]
The recurrence follows first for a finite head and then by absolute
convergence. Each \(E_{j,h}\) uses even moments through order \(2jh\).
Since \(B_j\) is the largest indexed product in \(E_{j,1}\),
\(B_j^h\le E_{j,h}\le B_j^{h-1}E_{j,1}\).
For \(h>1\) define
\(\underline B_{j,h}=(E_{j,h}/E_{j,1})^{1/(h-1)}\) and
\(\overline B_{j,h}=E_{j,h}^{1/h}\), with both equal to one at \(j=0\).
Then
\begin{equation}\label{eq:factor-interval}
 \sqrt{\frac{\underline B_{k-1,h}}{\overline B_{k,h}}}
 \le x_k\le
 \sqrt{\frac{\overline B_{k-1,h}}{\underline B_{k,h}}}.
\end{equation}
In particular \(x_k=\lim_{h\to\infty}(E_{k-1,h}/E_{k,h})^{1/(2h)}\)
for each fixed rank, including multiplicities. For finite input
intervals the alternating recurrence requires certified coefficient
errors before dividing by positive lower endpoints.

If the preceding factors are known, a separate estimate uses
\(s_h^{(k)}=p_h-\sum_{i<k}b_i^h\). The ratio
\(s_{h+1}^{(k)}/s_h^{(k)}\) is a weighted average of the factors
from index \(k\) onward. Subtract the average divided by \(b_k\)
from one, then use \(s_h^{(k)}\ge b_k^h\) and bound each
\((b_i/b_k)^{h-1}\) for \(i>k\) by the first such term. This gives
\[
 0\le 1-\frac{s_{h+1}^{(k)}}{b_k s_h^{(k)}}
 \le (b_{k+1}/b_k)^{h-1}\sum_{i>k} b_i/b_k .
\]
Strict separation gives geometric decay after the tail is bounded.
The subtraction also requires precision on the preceding factors
\cite[\S3]{Founding07}.

\paragraph{Exchanges across one rank.}
For \(b_{j+1}\le\theta_j\le b_j\), write
\[
 A_j=\sum_{i\le j}(\theta_j/b_i)^h,qquad
 C_j=\sum_{i>j}(b_i/\theta_j)^h .
\]
A term in \(E_{j,h}/B_j^h\) removes and inserts the same number of
indices across the cut at \(j\). All powers of \(\theta_j\)
therefore cancel. Writing \(e_r(\cdot)\) for an elementary symmetric
function here gives
\begin{equation}\label{eq:exchange-bound}
\begin{aligned}
 1\le \frac{E_{j,h}}{B_j^h}
 &=\sum_{r=0}^j e_r(((\theta_j/b_i)^h)_{i\le j})
                  e_r(((b_i/\theta_j)^h)_{i>j})\\
 &\le\sum_{r=0}^{\infty}\frac{(A_jC_j)^r}{(r!)^2}
 \le e^{A_jC_j}.
\end{aligned}
\end{equation}
The symmetric sums are bounded by the indicated powers divided by
\(r!\) on each side, by expanding the powers and counting the terms
with distinct indices. The tail sums converge by \(\sum_i b_i<\infty\).

If \(D_j\ge A_j C_j\) and \(0<L_j\le E_{j,h}\le U_j\), then
\begin{equation}\label{eq:exchange-interval}
 \left(\frac{L_{k-1}e^{-D_{k-1}}}{U_k}\right)^{1/(2h)}
 \le x_k\le
 \left(\frac{U_{k-1}e^{D_k}}{L_k}\right)^{1/(2h)} .
\end{equation}
Use \(D_0=0\) and \(L_0=U_0=1\).
Indeed, \(B_j^h\) lies between \(L_j e^{-D_j}\) and \(U_j\), and
\(x_k^{2h}=B_{k-1}^h/B_k^h\). The logarithm of the ratio of its two
endpoints is
\[
 \frac{D_{k-1}+D_k+\log(U_k/L_k)+\log(U_{k-1}/L_{k-1})}{2h}.
\]
A finite counted head through \(J\ge j+1\) supplies \(D_j\) if a
threshold can be chosen between the upper bound for \(b_{j+1}\) and
the lower bound for \(b_j\). Bound the known weights by interval
endpoints. For the rest, if \(b_{J+1}\le\beta\le\theta_j\) and
\(\sum_{i>J}b_i\le S\), the remaining contribution to \(C_j\)
is at most \((\beta/\theta_j)^{h-1}S/\theta_j\). The upper
endpoint for \(b_J\) can supply \(\beta\), and
\(\sum_i b_i=e_1\) supplies \(S\) after subtracting a lower
bound on the head sum. If both nearest threshold ratios are at
most \(q<1\), then \(A_j\le jq^h\) and
\(C_j\le q^{h-1}\sum_{i>j}b_i/\theta_j\). This gives geometric
decay for fixed \(j\). At a repeated factor \eqref{eq:factor-interval}
still applies.

At high ranks the spacing gives a uniform allowance for each fixed
\(T\). Take \(\theta_j=(b_j b_{j+1})^{1/2}\).
For \(\lfloor j/2\rfloor\le i\le 2j\), the adjacent differences
\(\log b_i-\log b_{i+1}=2\log(1+\Delta_i/x_i)\) are at least
\(c_T/j\) by \eqref{eq:high-rank-geometry}. The nearest weights on
each side of the cut thus cost half a step, and each subsequent
weight at least one additional step. On this range either sum is at
most \(e^{-c_T h/(2j)}/(1-e^{-c_T h/j})\).
On the far head, each base \(\theta_j/b_i\) is at most a fixed
\(q_0<1\) since \(x_{\lfloor j/2\rfloor}/x_j\to1/2\).
On the far tail each base \(b_i/\theta_j\) is also at most
\(q_0<1\) after increasing \(q_0\), since \(x_{2j}/x_j\to2\).
The head costs at most \(j q_0^h\), and the tail costs at most
\(q_0^{h-1}\theta_j^{-1}\sum_{i>2j}b_i=O_T(jq_0^{h-1})\).
For \(h\ge j\) these far bounds and both near sums are
\(O_T(e^{-c'_T h/j})\) with some \(c'_T>0\). For example,
one can absorb the factor \(j\) into half the exponential decay
in \(h\). We can therefore take \(D_j=C_Te^{-c_T'' h/j}\)
in \eqref{eq:exchange-interval} for sufficiently large \(j\)
and all \(h\ge j\).

\paragraph{Uncertainty in the symmetric coefficients.}
For \(\omega=e^{2\pi\ii/h}\), multiplying the rotated products gives
\[
 \prod_{l=0}^{h-1}G(\omega^l v)=\prod_i(1-b_i^h v^h),
 \qquad
 (-1)^j E_{j,h}=[v^{jh}]\prod_{l=0}^{h-1}G(\omega^l v).
\]
The product is locally absolutely convergent, and for \(j\le k\)
only the coefficients through \(kh\) are used. Equip these
truncated polynomials with
\(\|P\|_R=\sum_m |[v^m]P|R^m\), for \(R>0\).
The norm is submultiplicative and rotations are isometries.
Let \(\widetilde G\) be the center polynomial and \(d_R\) the
weighted sum of its coefficient radii in this norm. Put
\(U_R=\|\widetilde G\|_R+d_R\). Replacing the rotated factors
one at a time gives a telescoping sum. Taking its norm and then
extracting degree \(jh\) gives
\begin{equation}\label{eq:exchange-coefficient-error}
 |\widetilde E_{j,h}-E_{j,h}|
 \le h U_R^{h-1}d_R/R^{jh}.
\end{equation}
Enclose any arithmetic error separately.

Write \(M_R=G(-R)=F(\ii\sqrt R)\). The exact truncated norm is
at most \(M_R\). If \(d_R\le\eta M_R\), then
\(U_R\le(1+2\eta)M_R\). Since \(E_{j,h}\ge B_j^h\),
the relative error in \eqref{eq:exchange-coefficient-error} is at most
\[
 h\eta(1+2\eta)^{h-1}
 \left(\frac{M_R}{R^j B_j}\right)^h.
\]
Take \(R=(4\pi k/\log k)^2\). For \(j=k-1,k\),
\begin{equation}\label{eq:exchange-coefficient-growth}
 \frac{M_R}{R^j B_j}
 =\prod_{i\le j}(1+x_i^2/R)\prod_{i>j}(1+R/x_i^2)
 \le e^{C_T k}.
\end{equation}
Every head factor is bounded by a constant since \(x_j^2\asymp R\).
The tail logarithm is at most \(R\sum_{i>j}b_i=O_T(k)\).
Thus coefficient radii \(2^{-p} e_m\) through degree \(kh\)
are sufficient for relative output errors \(O(\tau)\) if
\(p\ge C_T'kh+\log_2(h/\tau)+C_T''\), with constants large enough
to absorb the factor \((1+2\eta)^{h-1}\) and the passage from
radii to endpoints.

\begin{theorem}[Sufficient product input budget]\label{thm:high-rank-exchange}
For each fixed positive all-real \(T\), at sufficiently large rank
\(k\), \(h\ge k\) and relative coefficient enclosure radii at most
\(\tau\le1/4\) for \(E_{k-1,h},E_{k,h}\), the width of
\eqref{eq:exchange-interval} in units of \(\Delta_k\) is
\[
 O_T\!\left(\frac{k}{h}(e^{-c_T h/k}+\tau)\right).
\]
In particular, \(h\asymp_T k(1+\log(1/\varepsilon))\) with a
sufficiently large constant and \(\tau\asymp\varepsilon\)
give width \(O_T(\varepsilon\Delta_k)\). The sufficient even
moment order is \(2kh\), and the sufficient input precision per
normalized coefficient is \(O_T(kh+\log(h/\tau))\) bits.
\end{theorem}
\begin{proof}
Substitute the high-rank bound for \(D_j\) into the logarithmic
width of \eqref{eq:exchange-interval}. Each log coefficient
ratio is \(O(\tau)\). The logarithmic width is \(O_T(1/h)\),
so exponentiating and using \(x_k/\Delta_k\asymp k\) proves
the width estimate. The precision bound follows from
\eqref{eq:exchange-coefficient-error}--\eqref{eq:exchange-coefficient-growth}.
\end{proof}
The sufficient asymptotic constants come from the spacing theorem.
At a specified rank one can use the certified head and tail
allowance above. The symmetric coefficients themselves are evaluated
directly from the moments.


\subsection{Direct evaluation from the moments}

Use the polynomial \(P_N(z)=\sum_{m=0}^N(-1)^m e_m z^{2m}\).
For real \(x\), Taylor's theorem through order \(2N+1\) inside the
positive cosine integral gives
\begin{equation}\label{eq:moment-real-error}
 |F(x)-P_N(x)|\le e_{N+1}|x|^{2N+2}.
\end{equation}
If \(\widetilde P_N\) uses coefficient centers with radii \(d_m\), add
\(\sum_{m=0}^N d_m X^{2m}\) on \(|x|\le X\) and use an upper
endpoint for \(e_{N+1}\). A numerical evaluation must also enclose its
arithmetic error.

The defining source gives refinable coefficient intervals at any fixed
\(T\). For example, let \(J_m=\int_0^\infty u^{2m}e^{Tu^2}\Phi(u)\dd u\).
To truncate at \(U_0\ge1\), set
\[
 p_m(u)=Tu^2+9u+2m\log u-\pi e^{4u},\qquad
 A=-p_m'(U_0)>0,qquad 16\pi e^{4U_0}\ge2T .
\]
These conditions hold for all sufficiently large \(U_0\).
Concavity on \([U_0,\infty)\) gives
\(p_m(u)\le p_m(U_0)-A(u-U_0)\). Hence the omitted integral is
at most \(2\pi^2 C_\Phi e^{p_m(U_0)}/A\), where
\(C_\Phi=\sum_{n\ge1}n^4 e^{-\pi(n^2-1)}\).
On \([0,U_0]\), truncating the square modes at \(n=M\) costs at most
\[
 2\pi^2 U_0^{2m+1} e^{TU_0^2+9U_0}
       \sum_{n>M}n^4 e^{-\pi(n^2-1)} .
\]
The displayed sums admit geometric upper bounds, for instance by
\(n^2-1\ge n-1\) and the power series for \(\sum n^4 q^{n-1}\) at
\(q=e^{-\pi}\). On the remaining compact integral, upper and lower
interval sums converge as the mesh tends to zero since the finite
integrand is continuously differentiable. Divide the resulting
enclosures as \(e_m=J_m/((2m)!J_0)\) using a positive lower endpoint
for \(J_0\). This gives a finite error allowance for generation of the
inputs. The order and precision estimates below concern the moments
supplied to the polynomial.

\begin{lemma}[Field bounds on circles near a high zero]\label{lem:high-rank-circle}
For every \(r\ge0\), \(\min_{|z|=r}|F(z)|=|F(r)|\).
There is a fixed \(d>0\) such that, for all sufficiently large
\(k\) and each fixed \(A>0\),
\begin{equation}\label{eq:direct-growth}
\begin{aligned}
 |F(r)|&\ge e^{-C_T k}\frac{|r-x_k|}{k\Delta_k}
             && (|r-x_k|\le d\Delta_k),\\
 F(2\ii X)&\le e^{C_{T,A}k}
             && (0\le X\le A k/\log k).
\end{aligned}
\end{equation}
\end{lemma}
\begin{proof}
On \(|z|=r\), the modulus of each factor in \eqref{eq:all-real-product}
is at least \(|1-b_i r^2|\), with equality for real \(z\). This
passes to the full product. The factor at \(i=k\) has magnitude
\(|r-x_k|(r+x_k)/x_k^2\gg_T |r-x_k|/(k\Delta_k)\).
For \(\lfloor k/2\rfloor\le i\le2k\), the gaps are uniformly at
least a fixed multiple of \(\Delta_k\) by \eqref{eq:high-rank-geometry}.
Take \(d\) below half this multiple. Then for \(i\ne k\),
\(|r-x_i|\gg_T |i-k|\Delta_k\) and
\(|1-r^2/x_i^2|\ge c_T |i-k|/k\), using \(x_i\ll_T k\Delta_k\)
and \(r+x_i\ge x_i\). The product of these bounds is
\(\ge e^{-C_T k}\) by \(m!\ge(m/e)^m\) on the two sides of \(k\).
Below \(\lfloor k/2\rfloor\) each factor has magnitude at least one
eventually, since \(x_{\lfloor k/2\rfloor}/x_k\to1/2\).
Above \(2k\), the quantities \(r^2/x_i^2\) are bounded away
from one since \(x_{2k}/x_k\to2\). Their sum is \(O_T(k)\) by
\eqref{eq:high-rank-geometry}, so the logarithm of the tail product
is bounded below by \(-C_T k\).

For the last assertion, split
\(\log F(2\ii X)=\sum_i\log(1+4X^2/x_i^2)\) at \(2k\).
The head satisfies \(X/x_i\le C_{T,A}k/i\) there by
\eqref{eq:high-rank-geometry}, with the comparison constant enlarged
to cover the finite prefix. For \(i\le2k\), \(1+4X^2/x_i^2\le
C_{T,A}'(2k/i)^2\). The head logarithm is \(O_{T,A}(k)\) by the
factorial bound. The tail logarithm is at most
\(4X^2\sum_{i>2k} b_i=O_{T,A}(k)\).
\end{proof}

\paragraph{Finite polynomial count.}
On \(|z|\le X\) a bound on the exact omitted tail is
\(4^{-(N+1)}F(2\ii X)\), obtained by extracting \(4^{-(N+1)}\)
from the positive coefficient sum at \(2\ii X\). There is also an
a posteriori test using a single extra coefficient. The \(e_m\)
are log-concave. To see this, for a finite product
\(f(v)=\prod_i(1+b_i v)\), let \(p=f^{(m-1)}\). Rolle's theorem
and the real roots of \(p\) give \((p')^2-pp''\ge0\), by
differentiating \(p'/p\) away from the roots and then using continuity.
At zero this gives \((m!e_m)^2\ge (m-1)!(m+1)!e_{m-1}e_{m+1}\)
for the finite coefficients. Pass to the locally uniform product limit.
For \(N\ge1\), the decreasing ratios yield
\begin{equation}\label{eq:moment-circle-error}
 \sup_{|z|\le X}|F(z)-\widetilde P_N(z)|
 \le \sum_{m=0}^N d_m X^{2m}
       +\frac{\overline e_{N+1}X^{2N+2}}{1-\bar\eta},
 \qquad
 \bar\eta=\frac{X^2\overline e_{N+1}}{\underline e_N}<1 ,
\end{equation}
where the lower coefficient endpoint must be positive. To justify a
strict comparison at a circle of radius \(r\le X\), bound
\(\min_{|z|=r}|\widetilde P_N(z)|\) below by a value larger than
\eqref{eq:moment-circle-error}. This can be certified by covering the
circle with sufficiently small arcs and bounding variation using the
polynomial derivative. Rouch\'e transfers its root count to \(F\).
The polynomial count can be determined by finite winding arithmetic
on the same circle. Its value is twice the number of positive zeros
inside, with multiplicity, since the full divisor is real and even.
Counts \(2(k-1)\) and \(2k\) on two radii place the simple \(x_k\)
strictly between them.

\begin{theorem}[Sufficient direct input budget]\label{thm:high-rank-direct}
At each fixed positive all-real \(T\), for sufficiently large \(k\) and
\(0<\varepsilon\le1/4\), the finite polynomial test can bracket
\(x_k\) to width \(O_T(\varepsilon\Delta_k)\) using
\(O_T(k+\log(1/\varepsilon))\) as both the maximum even moment order
and a sufficient relative precision in bits per normalized coefficient.
\end{theorem}
\begin{proof}
Choose \(X=8\pi k/\log k\) so that \(x_k\) is interior to the search
interval \([X/4,X]\) for all sufficiently large \(k\). A grid
of mesh \(c\varepsilon/\log k\) with small fixed \(c>0\) has points
on both sides of \(x_k\) between one and three mesh widths away
and within \(d\Delta_k\). At each of those radii the minimum
modulus of \(F\) is at least \(e^{-C_T k}\varepsilon/(C_T k)\)
by Lemma~\ref{lem:high-rank-circle}. The two exact counts there
are \(2(k-1)\) and \(2k\).

The true ratio in \eqref{eq:moment-circle-error} satisfies
\(X^2 e_{N+1}/e_N\le(e_N X^{2N})^{1/N}
\le 4^{-1}\exp(C_T k/N)\), because the earlier ratios have
product \(e_N X^{2N}\) and \(e_0=1\). For \(N\ge C'_T k\)
with sufficiently large \(C'_T\), small relative coefficient radii
keep \(\bar\eta\) bounded strictly below one. The resulting
geometric tail bound is \(O(4^{-(N+1)}e^{C_T k})\).
Uniform relative radii \(\delta\) in the lower coefficients
contribute at most \(\delta F(\ii X)\le\delta e^{C_T k}\).
Taking \(N\) and \(\log(1/\delta)\) each a sufficiently large
multiple of \(k+\log(1/\varepsilon)\) makes the sum smaller than a
quarter of the circle margin. The polynomial minimum therefore
exceeds the sum at both chosen radii, and sufficiently refined
interval arithmetic certifies these strict tests. The input order
is \(2N+2\).
\end{proof}

The finite test also applies at a specified lower rank without a
spacing theorem. Enumerate rational radii, moment orders and
precisions, interleaving the boundary checks. For a fixed radius
avoiding the divisor, compact convergence gives a positive margin
eventually. The ratio bound eventually succeeds because the series
is entire with positive coefficients and decreasing ratios.
Thus every specified simple root can be bracketed to arbitrarily
small width. Two such adjacent brackets give a positive lower
bound on their separation and can be refined to a prescribed
fraction of that bound. The theorem gives an asymptotic input budget
and leaves search, polynomial arithmetic and moment generation costs
to be quantified.

\subsection{Matching the absolute clock rank at \texorpdfstring{\(t=1/2\)}{t=1/2}}

For \(0<t\le1/2\), Polymath's Theorem~1.3 bounds the error of
the normalized approximation to \(H(t,x+\ii y)\) for \(x\ge200\),
\(0\le y\le1\) \cite{Polymath2019}. Its high-zero theorem uses
\[
 G_t(x)=\frac{x}{4\pi}\left(\log\frac{x}{4\pi}-1\right)
           +\frac{11}{8}+\frac{t}{16}\log\frac{x}{4\pi},
 \qquad G_t(u_j)=j .
\]
Theorem~1.5(i)--(ii) gives one simple real zero \(r_j\) in each
disk of radius \(c/\log u_j\) for \(j\ge\exp(C/t)\). Every zero
with real part \(x\ge\exp(C/t)\) is real and is within
\(O(x^{-ct})\) of a clock center.
Here \(C\) is a sufficiently large absolute constant and
\(c>0\) a sufficiently small one.
The \href{https://arxiv.org/html/1904.12438v2\#S9.E104}{proof after (104)} gives
\begin{equation}\label{eq:clock-error}
 |r_j-u_j|\ll\frac{u_j^{-ct}}{\log u_j},\qquad
 s_j=u_{j+1}-u_j\sim\frac{4\pi}{\log(u_j/(4\pi))}.
\end{equation}
The spacing follows by differentiation of \(G_t\). Theorem~1.5(iii)
gives \(N_t(X)=G_t(X)+O(1)\), where \(N_t\) counts zeros from the
imaginary axis by real part, with multiplicity. At an all-real
time this count is the positive rank from the origin
\cite[Theorem~1.5 and \S9]{Polymath2019}.

\begin{proposition}[Absolute clock rank]\label{prop:clock-rank}
At \(t=1/2\), \(r_j=x_j\) for all sufficiently large \(j\).
\end{proposition}
\begin{proof}
The divisor at \(t=1/2\) is all real by de Bruijn's upper bound
\cite[\S1]{Polymath2019}. The windows in \eqref{eq:clock-error}
are disjoint and ordered for large \(j\). Consequently for
\(x\ge b\), with \(b\) fixed sufficiently large and in a gap
between the windows, one has
\[
 N_t(x)=\lfloor G_t(x)\rfloor+d_0+\rho(x)
\]
for an integer \(d_0\). The error satisfies \(|\rho|\le1\), and
apart from interval endpoints its support lies between each clock
center and the corresponding root. Each dyadic interval \([R,2R]\)
contains \(O(R\log R)\) clock centers. By \eqref{eq:clock-error}
the integral of \(|\rho(x)|/x\) on that interval is
\(O(R^{-ct})\). Thus \(\rho\) is integrable against \(\dd x/x\)
on \([b,\infty)\).

On the imaginary axis use the complete product at time \(t\).
For \(Y>0\), integrate
\(K_Y(x)=2Y^2/[x(x^2+Y^2)]\) above each positive root.
The integral for root \(x_i\) is \(\log(1+Y^2/x_i^2)\).
Positivity permits summation under the integral, giving
\[
 \log\frac{H(t,\ii Y)}{Z(t)}
   =\int_0^\infty K_Y(x)N_t(x)\dd x .
\]
The finite head contributes \(O(1)\) as \(Y\to\infty\), since
\(K_Y(x)\le2/x\) and there is no zero near the origin.
The error \(\rho\) also costs \(O(1)\).
The function \(p(s)=\lfloor s\rfloor-s+1/2\) has a bounded
primitive by its zero mean on each period. Since \(G_t'\)
is eventually positive and \(G_t''(x)=1/(4\pi x)-t/(16x^2)\),
we may choose \(b\) so that \(G_t'\) is increasing there.
With \(s=G_t(x)\), the weight for \(p(s)\)
is \(K_Y(x)/G_t'(x)\), decreasing to zero and bounded at
\(b\) uniformly in \(Y\). Integration by parts shows that
this term is \(O(1)\). It remains to integrate
\(G_t(x)+d_0-1/2\) over \([b,\infty)\).

The term in \(x(\log(x/(4\pi))-1)/(4\pi)\) can be extended
down to zero with an \(O(1)\) error. Substitute \(x=Yu\).
The integral of \((1+u^2)^{-1}\) is \(\pi/2\), and its
integral against \(\log u\) vanishes by \(u\mapsto1/u\).
The resulting contribution is
\(Y(\log(Y/(4\pi))-1)/4\). The constant \(7/8+d_0\)
is multiplied by \(\log(1+Y^2/b^2)=2\log Y+O(1)\).
For the remaining logarithmic term, replacing \(K_Y(x)\)
by \(2\mathbf 1_{x\le Y}/x\) gives \(\log^2(Y/(4\pi))+O(1)\)
as its integral against \(\log(x/(4\pi))\). The difference
uses the weight
\[
 W(u)=\frac2u\left(\frac{1}{1+u^2}-\mathbf 1_{u\le1}\right)
\]
in the scaled variable. Its integral is zero by inversion,
its absolute logarithmic moment is finite, and its integral
in absolute value on \([0,b/Y]\) is \(O(Y^{-2})\).
Thus the replacement costs \(O(1)\) even after multiplication
by \(\log(Y/(4\pi))+\log u\). Altogether,
\begin{equation}\label{eq:clock-count-growth}
 \log H(t,\ii Y)
 =\frac{Y}{4}\left(\log\frac{Y}{4\pi}-1\right)
   +\frac{t}{16}\log^2\frac{Y}{4\pi}
   +\left(\frac74+2d_0\right)\log Y+O(1).
\end{equation}

To evaluate the same logarithm from the source, put
\(F_0(Y)=H(0,\ii Y)\).
By evenness and the zero-time normalization,
\(F_0(Y)=\xi((1+Y)/2)/8\) for \(Y>0\). At
\(s=(1+Y)/2\), Stirling's formula and \(\zeta(s)=1+o(1)\)
give \((s/2)\log(s/(2\pi))-s/2+(3/2)\log s+O(1)\)
for \(\log F_0(Y)\). Substituting \(s\) gives
\[
 \log F_0(Y)=P(Y)+O(1),\qquad
 P(Y)=\frac{Y}{4}\left(\log\frac{Y}{4\pi}-1\right)
                       +\frac74\log Y .
\]
In particular the shift \(Y+1\) contributes \(\tfrac14\log Y\).
For a centered Gaussian \(V\) of variance \(2t\), the positive
source integral gives \(H(t,\ii Y)=\E F_0(Y+V)\).
Indeed the Gaussian average of \(\cosh((Y+V)u)\) is
\(e^{tu^2}\cosh(Yu)\), with interchange justified by positivity.

Write \(A_Y=P'(Y)=\tfrac14\log(Y/(4\pi))+7/(4Y)\).
On \(|v|\le Y^{1/4}\), Taylor's theorem gives
\(P(Y+v)=P(Y)+A_Yv+O(Y^{-1/2})\). Multiplying the
Gaussian density by \(e^{A_Yv}\) gives a shifted density of
mean \(2tA_Y=O(\log Y)\) times the factor \(e^{tA_Y^2}\).
The shifted Gaussian probability of the central interval tends
to one. On \(Y^{1/4}<|v|\le Y/2\), the bound
\(F_0(Y+v)\ll e^{P(Y)+C|v|\log Y}\) follows from the
estimate for \(P'\). The quadratic Gaussian exponent
dominates the linear term and gives a contribution of
\(O(e^{P(Y)-c_tY^{1/2}})\). On \(|v|>Y/2\),
evenness and the Stirling bound give
\(\log F_0(Y+v)=O((1+|v|)\log(2+|v|))\) since
\(|Y+v|\le3|v|\). This range contributes
\(O(e^{-c_tY^2})\). It follows that
\[
 \log H(t,\ii Y)=P(Y)+tA_Y^2+O(1)
               =P(Y)+\frac{t}{16}\log^2\frac{Y}{4\pi}+O(1).
\]
Comparison with \eqref{eq:clock-count-growth} forces \(d_0=0\).
The point with \(G_t(X)=j+1/2\) is eventually outside the
windows and has exactly \(j\) positive zeros through it.
This gives the asserted rank.
\end{proof}

The full published displacement at these matched ranks gives
position and gap errors \(O(u_k^{-ct})\) in units of \(s_k\).
For the gap, subtract the two displacements in
\eqref{eq:clock-error} and use \(u_{k+1}/u_k\to1\) and the
displayed spacing. This calibration leaves the height onset
unevaluated. A finite alternative at a cutoff in a proved clock
gap is an exact prefix count, for example by the finite
polynomial count above. If the adjacent clock labels there are
\(j,j+1\), subtract \(j\) from that count to obtain the offset.

For refinement with moments at \(t=1/2\), the matched clock
gives a short real interval containing just \(x_k\), for
example \([u_k-d's_k,u_k+d's_k]\) for small fixed \(d'>0\).
By \eqref{eq:clock-error} the whole interval eventually lies
within the neighborhood in Lemma~\ref{lem:high-rank-circle}.
Its endpoints have opposite signs and magnitudes
\(\gg e^{-Ck}/k\). With coefficient and arithmetic errors
small compared with \(e^{-Ck}\varepsilon/k\), the real error
bound \eqref{eq:moment-real-error} can be used to bisect this
interval. If a test interval contains zero, the modulus of the
true field is at most twice the certified radius, and the
lower bound in Lemma~\ref{lem:high-rank-circle} puts the test
point within \(O(\varepsilon s_k)\) of the root. Otherwise the
sign selects the next half-interval. The same moment order and
sufficient precision as in Theorem~\ref{thm:high-rank-direct}
meet this error allowance by \eqref{eq:direct-growth}.

With the same coarse geometry, Theorem~\ref{thm:high-rank-direct}
uses a shorter moment prefix at large rank than the sufficient
budget in Theorem~\ref{thm:high-rank-exchange}. If the full
clock displacement is available, it already supplies the
displayed localization accuracy at \(t=1/2\) without acquiring
moments. For an additional prescribed refinement the theorems
give sufficient orders and precision with either construction.
Subtracting the neighboring root intervals gives gap uncertainty
\(O_T(\varepsilon\Delta_k)\) since successive spacings are
asymptotically equal.
They leave the onset and arithmetic work unevaluated, as well as
the height overlap with the large-time inverse in
Section~\ref{sec:anchors} at a common \(T\).

\subsection{The published time-zero benchmark}

The identity \(H(0,z)=\xi((1+\ii z)/2)/8\) gives \(x=2\gamma\)
for a critical-line zero of ordinate \(\gamma\). Platt and
Trudgian certify that the first \(12\,363\,153\,437\,138\) zeros
of positive ordinate through \(\gamma=3\,000\,175\,332\,800\)
lie on the line, corresponding to \(x=6\,000\,350\,665\,600\)
\cite[Theorem~1 and \S2]{PlattTrudgian2021}. Their sign-change
count with Turing's method certifies completeness at time zero.
A positive-time count at a specified vertical boundary also
requires accounting for boundary crossings or counting at that time.

Polymath proved \(\Lambda\le0.22\)
\cite[Theorem~1.1]{Polymath2019}. The second row of its
Table~1 uses an ordinate verification above \(2.51\cdot10^{12}\),
so the larger Platt--Trudgian verification gives \(\Lambda\le0.2\)
by \href{https://arxiv.org/pdf/2004.09765v1\#page=5}{Corollary~2}
of \cite{PlattTrudgian2021}. The criterion giving
\(\Lambda\le t_0+y_0^2/2\) combines a time-zero verification
through ordinate \(X/2\), final-time nonvanishing for
\(x\ge X+\sqrt{1-y_0^2}\) and \(y_0\le y\le\sqrt{1-2t_0}\),
and an intermediate-time barrier on
\(X\le x\le X+\sqrt{1-y_0^2}\) above
\(\sqrt{y_0^2+2(t_0-t)}\) up to \(\sqrt{1-2t}\)
\cite[Theorem~1.2]{Polymath2019}. An improvement by this
criterion requires these conditions at parameters satisfying
\(t_0+y_0^2/2<0.2\). The factor and circle estimates concern
one independently established all-real time.
